% Evans second edition, section 2.4.
% English source: printed p. 63; PDF page 80 (1-based).
% Chinese section starts checked from chapter opening, PDF 31; mean-value pages checked at PDF 36-37 and 53-55.
% Status: editorial study draft with reconstructed proofs and mathematical annotations.
\section{Wave Equation}
\label{sec:2-4}

\subsection{Spherical means and explicit solutions (Evans \S2.4.1)}
The wave equation $u_{tt}-\Delta u=0$ needs two initial data:
$u(\cdot,0)=g$, $u_t(\cdot,0)=h$.
Unlike the heat equation, it scales with $t\mapsto\lambda t$ when
$x\mapsto\lambda x$.

\paragraph{One dimension: d'Alembert (source Theorem 1).}
Write $\xi=x+t$, $\eta=x-t$. Then
$\partial_t^2-\partial_x^2=-4\partial_\xi\partial_\eta$.
Thus $u_{\xi\eta}=0$, and integration gives $u=F(\xi)+G(\eta)$.
At $t=0$, $F+G=g$ and $F'-G'=h$, whereas $F'+G'=g'$.
Solving for $F',G'$ and integrating yields
\[
 u(x,t)=\frac{g(x+t)+g(x-t)}2
              +\frac12\int_{x-t}^{x+t}h(y)\,dy .
 \tag{DA}
\]
For $g\in C^2$, $h\in C^1$, direct differentiation verifies the equation,
$u(\cdot,0)=g$, and $u_t(\cdot,0)=h$.
Only initial points in $[x-t,x+t]$ matter.

\Needspace{45mm}
\PDEAnnotation{两个初值各管什么}{初始位移 $g$ 生成两列反向平移的形状；初始速度 $h$ 在区间内累积。热方程只需一个初值，而波动方程的二阶时间导数要求两个。}


\PDEBodyDiagram{2.4}{The wave equation has a finite domain of dependence.}
{For speed one in one dimension, the solution at $(x,t)$ uses initial data only on $[x-t,x+t]$. At an intermediate time $s$, the relevant interval is $|y-x|\leq t-s$. In higher dimensions this interval becomes a ball. A forced equation also accumulates sources from the past cone.}
{ch02-wave-cone}

\paragraph{The common spherical-mean identity.}
Fix a spatial center $x$ and set
$M(r,t)=|\partial B_r|^{-1}\int_{\partial B_r(x)}u(y,t)\,dS$.
The radius-flux computation used for harmonic means now gives
\[
 M_r=\frac1{n\omega_n r^{n-1}}\int_{B_r(x)}\Delta u,\qquad
 (r^{n-1}M_r)_r=r^{n-1}M_{tt}.
\]
Thus $M$ satisfies the Euler--Poisson--Darboux equation
\[
 M_{tt}=M_{rr}+\frac{n-1}rM_r,\qquad
 M(0,t)=u(x,t),\quad M_r(0,t)=0 .
 \tag{EPD}
\]
The second radius derivative differentiates the ball integral by shells;
the wave equation changes $\Delta u$ into $u_{tt}$.
\Needspace{45mm}
\PDEAnnotation{同一个通量公式三次出现}{Laplace 方程令球面平均对半径不变；波动方程把同一半径变化联系到二阶时间导数；热方程则需要带方向的时空平均。不要因都出现球面而混为一谈。}

\paragraph{Three dimensions: Kirchhoff.}
For $n=3$, $W(r,t)=rM(r,t)$ satisfies $W_{tt}=W_{rr}$.
Extend $M$ evenly in $r$, hence $W$ oddly. Its initial data are
$rS_g(x,r)$ and $rS_h(x,r)$.
Use (DA) for $W$ and differentiate in $r$ at zero.
Since $W_r(0,t)=u(x,t)$, this gives
\[
 u(x,t)=\partial_t\big[tS_g(x,t)\big]+tS_h(x,t).
 \tag{KIR}
\]
Equivalently
\[
 u(x,t)=\frac1{4\pi}\int_{\partial B_1}
 [g(x+t\theta)+tDg(x+t\theta)\cdot\theta+t h(x+t\theta)]\,dS_\theta.
\]
The latter fixed-sphere expression makes initial limits transparent.
The formula solves the wave equation: by (EPD) applied to spatial
spherical means of $g$ or $h$, $tS_h$ satisfies the wave equation
in the variables $(x,t)$; time differentiation handles the $g$ term.
At zero, $tS_h=t h(x)+O(t^3)$ and
$\partial_t(tS_g)=g(x)+O(t^2)$.
These observations prove existence as well as the derivation from a
hypothetical solution, for smooth data.

\paragraph{Two dimensions: descent from three dimensions.}
Make $g,h$ independent of a third coordinate.
Project the upper and lower hemispheres onto the disk. Each contributes
surface Jacobian $t/\sqrt{t^2-|y-x|^2}$.
Formula (KIR) becomes
\[
 u(x,t)=\partial_t\left[\frac1{2\pi}
          \int_{B_t(x)}\frac{g(y)}{\sqrt{t^2-|y-x|^2}}\,dy\right]
       +\frac1{2\pi}\int_{B_t(x)}
          \frac{h(y)}{\sqrt{t^2-|y-x|^2}}\,dy .
 \tag{W2}
\]
Use $y=x+tz$ before differentiating: the singularity
$(1-|z|^2)^{-1/2}$ is integrable on the fixed unit disk.
This justifies differentiation for sufficiently smooth data.
Descent preserves the PDE because the third-coordinate derivative is
zero; the three-dimensional initial limits give the desired two data.
\Needspace{45mm}
\PDEAnnotation{降维为何出现权重}{投影把球面压到圆盘，靠近圆周时球面近乎竖直，面积 Jacobian 变大，产生 $(t^2-|y-x|^2)^{-1/2}$。它是几何投影权重，不是热球的权重。}

\paragraph{All odd dimensions (source Theorem 2).}
For $n=2k+1\geq3$, put $D_r=r^{-1}\partial_r$ and
$\gamma_n=1\cdot3\cdots(n-2)$. Define
\[
 S_n(t)f=\frac1{\gamma_n}D_t^{\,k-1}
                         \big[t^{n-2}S_f(x,t)\big].
\]
Then
\[
 u(x,t)=\partial_t S_n(t)g+S_n(t)h.
 \tag{WO}
\]
Here $S_n(t)$ denotes a solution operator, whereas $S_f(x,t)$ is the
normalized spherical average defined earlier.

\begin{expandedproof}[the radial intertwining step]
Let $E=r\partial_r$ and
$A_kF=D_r^{k-1}(r^{2k-1}F)=rP(E)F$,
where $P(E)=\prod_{j=1}^{k-1}(E+2j+1)$; for $k=1$ the product is one.
This factorization follows successively from
$D_r(r^aF)=r^{a-2}(E+a)F$.
The radial Laplacian is $r^{-2}E(E+2k-1)$.
The polynomial identity
\[
 (E+1)P(E)=(E+2k-1)P(E-2)
\]
therefore proves
\[
 (A_kF)''=A_k\left(F''+\frac{2k}{r}F'\right).
\]
For $F(r)=S_f(x,r)$ the bracket equals $\Delta_xF$ by the divergence
identity behind (EPD). Hence $A_kS_f$, now with $r=t$, solves the wave
equation. It commutes with $\Delta_x$ since the operator acts only on $r$.

Spherical symmetry gives $S_f(x,t)=f(x)+O(t^2)$ with an even smooth
extension at zero. The coefficient of $t f(x)$ in
$A_kS_f$ is $P(0)=3\cdot5\cdots(2k-1)=\gamma_n$.
All other terms start at $t^3$.
Thus $S_n(t)f=t f+O(t^3)$, giving initial value zero and initial
velocity $f$. Its time derivative has initial value $f$ and initial
velocity zero. This proves (WO) for smooth data and explains its constants.
Finite differentiability follows by differentiating the fixed-sphere
integrals; the source sufficient assumptions are
$g\in C^{m+1}$, $h\in C^m$, $m=(n+1)/2$.
\end{expandedproof}

\paragraph{All even dimensions (source Theorem 3).}
For even $n\geq2$, put $\gamma_n=2\cdot4\cdots n$ and
\[
 B_f(x,t)=\frac{t^n}{\omega_nt^n}
          \int_{B_t(x)}\frac{f(y)}{\sqrt{t^2-|y-x|^2}}\,dy
 =\frac1{\omega_n}\int_{B_t(x)}
          \frac{f(y)}{\sqrt{t^2-|y-x|^2}}\,dy .
\]
Set $S_n(t)f=\gamma_n^{-1}D_t^{(n-2)/2}B_f(x,t)$.
The same formula $u=\partial_tS_n(t)g+S_n(t)h$ holds.
To derive it, extend the data constantly in one extra coordinate,
apply (WO) in dimension $n+1$, and project both hemispheres.
Their Jacobian gives $2t/\sqrt{t^2-|y-x|^2}$; the sphere and ball
area constants simplify to the displayed $\gamma_n$.
The inherited PDE and initial limits prove it.
For smooth data there is no endpoint issue; for the source finite
regularity statement take $m=(n+2)/2$, $g\in C^{m+1}$, $h\in C^m$.
When $n=2$, this reduces to (W2).

For odd $n\geq3$, only data on the sphere at distance $t$ enter the
formula, the sharp Huygens phenomenon.
For even $n$, the whole ball enters and a trailing influence remains.
Dimension one is an exception to the odd-dimensional sharp-front
statement: (DA) includes an interval integral of $h$.

\subsection{Forcing (Evans \S2.4.2, Theorem 4)}
The operator $S_n(t)$ above solves the homogeneous equation with
initial displacement zero and velocity $f$.
For smooth forcing set
\[
 v(t)=\int_0^t S_n(t-s)f(s)\,ds.
\]
Since $S_n(0)=0$ and $\partial_tS_n(0)=I$, differentiate twice:
\[
 v_t=\int_0^t\partial_tS_n(t-s)f(s)\,ds,\qquad
 v_{tt}=f(t)+\int_0^t\partial_t^2S_n(t-s)f(s)\,ds.
\]
The integrals in $v_{tt}$ and $\Delta v$ cancel.
Thus $v_{tt}-\Delta v=f$, with both initial data zero.
Add the homogeneous solution for general initial data.

\subsection{Energy and finite propagation (Evans \S2.4.3, Theorems 5--6)}
For a smooth solution with fixed zero Dirichlet boundary values,
\[
 \mathcal E(t)=\frac12\int_U(u_t^2+|Du|^2),\qquad
 \mathcal E'(t)=\int_{\partial U}u_t\partial_\nu u=0.
\]
The boundary term vanishes because $u=0$ there for all times implies
$u_t=0$ there. Apply this to the difference for uniqueness; zero initial
energy gives $Du=u_t=0$, and zero initial displacement removes constants.

For local propagation, consider the shrinking ball
$B_{R-t}(x_0)$ and the cone energy
\[
 e(t)=\frac12\int_{B_{R-t}(x_0)}(u_t^2+|Du|^2).
\]
Differentiating the moving domain contributes a negative surface term.
The PDE and spatial integration by parts yield
\[
 e'(t)=\int_{\partial B_{R-t}}
 \left[u_t\partial_\nu u-\tfrac12(u_t^2+|Du|^2)\right]\,dS\leq0.
\]
The inequality follows from
$2u_t\partial_\nu u\leq u_t^2+|\partial_\nu u|^2\leq u_t^2+|Du|^2$.
If both initial data vanish on $B_R(x_0)$, $e(0)=0$ and hence
$e(t)=0$ for $t<R$. Along each vertical segment inside the cone
$u_t=0$ and its initial value is zero, so $u=0$ throughout the cone.
Continuity gives the apex value when applicable.
\Needspace{45mm}
\PDEAnnotation{传播速度来自能量通量}{球半径以速度 $1$ 收缩时，移边界损失足以压住波的能量流入。热核虽处处为正，波的锥外影响却严格为零；这是两个方程的本质区别。}
